\documentclass[../../main.tex]{subfiles}
\begin{document}

{\bf [解]}

\begin{figure}[htb]
  \centering
  \begin{tikzpicture}[scale=1.4, >=stealth, every node/.style={font=\small}]

    % --- 1. 三角形の各頂点の定義 ---
    % B を原点、BC を x 軸上にとる (a = 5.0, c = 3.6, B = 55°)
    \coordinate (B) at (0, 0);
    \coordinate (C) at (5.0, 0);
    \coordinate (A) at (55:3.6);

    % --- 2. 辺 BC の n 等分点 P_k の配置 ---
    % n = 7 と見立てた分割例
    \coordinate (P1) at ($ (B)!1/7!(C) $);
    \coordinate (P2) at ($ (B)!2/7!(C) $);
    \coordinate (Pk) at ($ (B)!4/7!(C) $); % 代表点 P_k (k=4)
    \coordinate (Pk1) at ($ (B)!5/7!(C) $); % P_{k+1}

    % --- 3. 角 B の円弧 ---
    \draw[thick, purple!80!black] (0.55, 0) arc (0:55:0.55);
    \node[purple!80!black, font=\footnotesize] at (27.5:0.8) {$B$};

    % --- 4. 三角形 ABC の枠線 ---
    \draw[very thick, black] (A) -- (B) -- (C) -- cycle;

    % 辺の長さラベル
    \node[left=2pt] at ($ (A)!0.5!(B) $) {$c$};
    \node[above right=2pt] at ($ (A)!0.5!(C) $) {$b$};

    % --- 5. 線分 AP_k および補助線 ---
    % A から P_k への線分（主役）
    \draw[very thick, red!85!black] (A) -- (Pk);
    \node[red!85!black, font=\footnotesize, above right=-1pt] at ($ (A)!0.48!(Pk) $) {$\mathrm{AP}_k$};

    % その他の分割点への細い点線
    \draw[dashed, gray!60] (A) -- (P1);
    \draw[dashed, gray!60] (A) -- (P2);
    \draw[dashed, gray!60] (A) -- (Pk1);

    % --- 6. 各点のプロットとラベル ---
    \fill (A) circle (1.2pt) node[above=2pt] {$\mathrm{A}$};
    \fill (B) circle (1.2pt) node[below left=2pt] {$\mathrm{B} = \mathrm{P}_0$};
    \fill (C) circle (1.2pt) node[below right=2pt] {$\mathrm{C} = \mathrm{P}_n$};

    \fill[gray!80!black] (P1) circle (1.0pt) node[below=2pt, font=\scriptsize] {$\mathrm{P}_1$};
    \fill[gray!80!black] (P2) circle (1.0pt) node[below=2pt, font=\scriptsize] {$\mathrm{P}_2$};
    \node[below=2pt, font=\scriptsize, gray!80!black] at ($ (P2)!0.5!(Pk) $) {$\dots$};

    \fill[red!85!black] (Pk) circle (1.3pt) node[below=2pt, font=\footnotesize, red!85!black] {$\mathrm{P}_k$};
    \fill[gray!80!black] (Pk1) circle (1.0pt) node[below=2pt, font=\scriptsize] {$\mathrm{P}_{k+1}$};
    \node[below=2pt, font=\scriptsize, gray!80!black] at ($ (Pk1)!0.5!(C) $) {$\dots$};

    % --- 7. 長さの寸法線 (BC = a, BP_k = ak/n, P_k P_{k+1} = a/n) ---
    % 1区間の長さ P_k P_{k+1}
    \draw[<->, teal!85!black, thick] ($ (Pk) + (0, 0.22) $) -- ($ (Pk1) + (0, 0.22) $)
    node[midway, above=1pt, font=\scriptsize] {$\frac{a}{n}$};

    % BP_k の長さ
    \draw[<->, blue!85!black] ($ (B) + (0, -0.6) $) -- ($ (Pk) + (0, -0.6) $)
    node[midway, fill=white, inner sep=1pt, font=\scriptsize] {$\mathrm{BP}_k = \frac{ak}{n}$};
    \draw[dotted, gray!70] (Pk) -- ($ (Pk) + (0, -0.65) $);

    % 辺 BC 全体の長さ a
    \draw[<->, gray!80!black] ($ (B) + (0, -1.05) $) -- ($ (C) + (0, -1.05) $)
    node[midway, fill=white, inner sep=1pt, font=\footnotesize] {$\mathrm{BC} = a$};
    \draw[dotted, gray!70] (B) -- ($ (B) + (0, -1.1) $);
    \draw[dotted, gray!70] (C) -- ($ (C) + (0, -1.1) $);

  \end{tikzpicture}
  \caption{辺 $\mathrm{BC}$ の $n$ 等分点 $\mathrm{P}_k$ と $\triangle \mathrm{ABP}_k$ における幾何配置}
  \label{fig:triangle_division_APk}
\end{figure}

題意から$k=1,2,\cdots,n-1$に対して
\begin{align}
  \mathrm{P_kP_{k+1}}=\dfrac{a}{n} \label{eq:1}
\end{align}
である．
三角形$\mathrm{ABP_k}$に余弦定理を用いて
\begin{align}
  \mathrm{AP_k}^2
  &=\mathrm{AB}^2+\mathrm{BP_k}^2-2\mathrm{AB}\cdot\mathrm{BP_k}\cos \angle\mathrm{B} \\
  &=c^2+\left(\dfrac{ak}{n}\right)^2-2c\dfrac{ak}{n}\cos\angle\mathrm{B} \label{eq:2}
\end{align}
となる．
また，三角形ABCに余弦定理を用いて
\begin{align}
  \mathrm{AC}^2
  &=\mathrm{AB}^2+\mathrm{BC}^2-2\mathrm{AB}\cdot\mathrm{BC}\cos \angle\mathrm{B} \\
  \therefore
  b^2&=c^2+a^2-2ac\cos\angle\mathrm{B} \label{eq:3}
\end{align}
となる．

\cref{eq:2,eq:3}から$\cos\angle\mathrm{B}$を消去して
\begin{align}
  \mathrm{AP_k}^2=c^2+\left(\dfrac{ak}{n}\right)^2-(a^2+c^2-b^2)\dfrac{k}{n} \label{eq:4}
\end{align}
を得る．
$A=a^2+c^2-b^2$とおいて，\cref{eq:4}を$k=1,2,\cdots,n$について足して
\begin{align}
  \dfrac{1}{n}\sum_{k=1}^n\mathrm{AP_k}^2
  =c^2+\dfrac{1}{n}\sum_{k=1}^n\left(\left(\dfrac{ak}{n}\right)^2-A\dfrac{k}{n}\right)
\end{align}
であるから，$n\to\infty$の極限をとって区分求積法を用いて
\begin{align}
  \lim_{n\to\infty}\dfrac{1}{n}\sum_{k=1}^n\mathrm{AP_k}^2&= %
  c^2+\int_0^1(a^2x^2-Ax)\, dx \\
  &=c^2+\left[\dfrac{a^2}{3}x^3-\dfrac{1}{2}Ax^2\right]_{0}^{1} \\
  &=c^2+\dfrac{a^2}{3}-\dfrac{1}{2}(a^2+c^2-b^2) \\
  &=\dfrac{1}{6}(-a^2+3b^2+3c^2)
\end{align}
となる．$\cdots$(答)
\newpage
\end{document}
